% The details of the mechanical checks, moved here from Section 3 on 2026-10-14.
\section{The mechanical checks in detail}
\label{sec:checks}

\paragraph{The checks.} Before a search, a run of spaces and line breaks becomes one space on both sides, and the Markdown marks at the start of a line are taken off. Punctuation, digits and word order have to match. The checks follow in the order of the report, each with the heading of its findings.
\begin{points}
\item \emph{Every source segment is in the merge word for word, or has a record.} A search for the segment's text. A segment that is not found and has no record is reported under one of three headings: ``Reworded and undeclared'' if the program found it in altered form, ``Absent and undeclared'' if it did not, which is the silent loss, and ``Heading dropped in silence'' for a heading that the merge did not take while its section is still there. The paragraphs after this list say how the program tells them apart. A segment that is found although a record says it was subsumed or superseded is reported as ``Declared gone, still here''.
\item \emph{Every record names a real segment.} A lookup among the source segments. A record for a segment that no source has is an ``Invented segment''.
\item \emph{Every replacement is in the merge.} A search for the replacement's text. If it is not there, the finding is ``Unresolved replacement''.
\item \emph{Every record is of a kind its fidelity level allows.} A lookup in a table of levels and kinds. A kind the level does not allow is ``Not permitted here''.
\item \emph{Every protected value is in the merge unchanged.} Patterns find the values in a source segment, and a search looks for each of them in the merge. A value that is not found is a ``Verbatim violation''.
\item \emph{The records keep within the two limits.} The program counts declared drops against the loss budget of \ccnum{tool.loss_budget_percent}\%, and missing segments per replacement against the limit of \ccnum{tool.fan_in}. Either limit exceeded is ``Over budget''.
\item \emph{The merge has a title.} It needs one if a source has one, and under the stricter title settings it must be a source's title, character for character. Otherwise the finding is ``Title not from a source''.
\item \emph{A source title that was not taken has a record.} The record must say superseded, and its replacement must contain the merge's title. Without such a record the finding is ``Title dropped in silence''.
\item \emph{The merge says nothing twice.} Its segments are compared with each other, titles and headings left out. Two segments that are equal word for word are reported as ``Stated twice''.
\end{points}

\paragraph{How a segment is looked for.} A source segment that is not in the merge word for word is looked for in two further ways. The first is the similarity score. It is the ratio from \texttt{difflib}: twice the number of matching characters, divided by the length of both texts together. The segment is scored against every segment of the merge, and the best score counts. At or above the cut-off of \ccnum{tool.near_match} the segment counts as reworded. The score reaches the cut-off when a little more than half of a sentence is left, so a sentence that was cut short can pass as reworded. It is still reported, with the difference printed, but not as a silent loss (Section~\ref{sec:limitations}).

The second is for a sentence, a list item or a table row that the score left below the cut-off. The program asks whether the words of the segment still stand in the merge in their order, as one unbroken run, with punctuation and capitals ignored. That finds a sentence which the merge kept and joined to its neighbour: the word-for-word search misses it because its full stop is gone, and the score misses it because the joined sentence is much longer. The run may also reach across two neighbouring segments of the merge, for a sentence that was split in two or run together with the next. There a short list of connecting words may differ: the articles, ``and'', the forms of ``to be'', and pronouns such as ``it'' and ``which''. No other word may be missing or added, and every number, link and code value of the segment must be there. A segment that is found in either way is reworded. One that is found in neither is missing.

A heading is looked for only among the headings of the merge. A source heading that is not there and has no record is a silent loss if the whole section under it is missing too, or if it carries a number, a link or a code value that the merge does not have. Otherwise its section is in the merge under another heading, and only the label was passed over. The program reports that as ``Heading dropped in silence''. I call it an \emph{untaken heading}. It is a finding about the merge's records and not about its content, so the run ends with exit code 3.

\paragraph{The patterns for protected values.} The patterns find inline code, links, URLs, file paths, version strings and numbers. A number is a run of digits with what is attached to it, such as \texttt{30s} or \texttt{100\%}, and a number written as a word is not found. Each value is searched for whole, so \texttt{512} is not found inside \texttt{5120}. At the default level a missing value is excused in two cases. Its segment has a record of the kind dropped, which the loss budget counts. Or it has a record of the kind superseded whose replacement is in the merge and is the wording of another source segment. The second case is for sources that write one number in two ways, such as ``two billion'' and ``2,000,000,000''.
